\documentclass[conference]{IEEEtran}
\IEEEoverridecommandlockouts
\usepackage{cite}
\usepackage{amsmath,amssymb}
\usepackage{algorithm}
\usepackage{algorithmic}
\usepackage{graphicx}
\usepackage{booktabs}
\usepackage{multirow}
\usepackage{textcomp}
\usepackage{url}

\begin{document}

\title{Deadline-Aware Elastic Provisioning for\\Hibernation-Tolerant Bag-of-Tasks Scheduling\\on Spot and Burstable Cloud Instances}

\author{\IEEEauthorblockN{}\IEEEauthorblockA{}}

\maketitle

\begin{abstract}
Spot virtual machines cut cloud compute prices substantially but may be
hibernated by the provider at any time, which threatens deadline-constrained
workloads. Burst-HADS tolerates such interruptions by migrating the tasks of a
hibernated instance onto burstable VMs running in burst mode, falling back to
on-demand capacity when no burstable instance is free. Because a burstable
instance accepts only one task at a time and the number provisioned is a fixed
fraction of the selected fleet, the rescue capacity available at the moment of
an interruption is bounded independently of how much work is actually
displaced. We present R-BurstHADS, which extends Burst-HADS's dynamic scheduler
with elastic spot provisioning drawn from the same instance catalogue, sized by
a work-conservation argument against the remaining deadline rather than by a
fixed queue-length constant. Across 2{,}700 simulation runs spanning five
hibernation scenarios, six workload sizes and 30 seeds, R-BurstHADS reduces
makespan relative to Burst-HADS by 26.0--39.8\% in every scenario, while its
monetary-cost advantage scales monotonically with the interruption rate: it is
cost-neutral at $k_h=1$, 12.5\% cheaper at $k_h=3$, and 14.3--64.3\% cheaper at
$k_h=5$. All three schedulers complete 100\% of tasks, but only the two
burstable-aware schedulers meet the deadline in every configuration. We first
validate our HADS and Burst-HADS implementations against the published results
of the source work, reproducing the reported ordering and the direction and
magnitude of its cost differences.
\end{abstract}

\begin{IEEEkeywords}
cloud computing, spot instances, burstable instances, task scheduling,
hibernation, bag-of-tasks, deadline constraints
\end{IEEEkeywords}

\section{Introduction}

Infrastructure-as-a-Service providers sell surplus capacity at a steep discount
under names such as Amazon EC2 Spot. The discount is paid for in reliability:
the provider may reclaim the instance at short notice. Amazon EC2 additionally
offers \emph{hibernation}, in which the instance's memory and context are
written to its EBS root volume and execution is suspended rather than
destroyed, with the instance resuming when capacity becomes available again.
For a workload with a deadline, hibernation is not free even though the
checkpoint preserves progress: the capacity disappears for an unknown period,
and the displaced work must be placed somewhere else in time to finish.

Bag-of-Tasks (BoT) applications --- parameter sweeps, Monte Carlo simulation,
chromosome mapping, image processing --- are a natural fit for this market.
Their tasks are independent, so a lost instance costs only the work resident on
it, and the remaining tasks can in principle be redistributed freely.

Teylo et al. addressed this setting in two stages. HADS~\cite{hads2021}
schedules a BoT application across spot and non-burstable on-demand instances
and, when a spot VM hibernates, migrates its unfinished tasks to other running
instances or to newly launched on-demand ones. Burst-HADS~\cite{bursthads}
extends this with \emph{burstable} instances, which accumulate CPU credits
while idle and can then be driven at full processing power in burst mode; a
hibernation is answered first by migrating a task to an idle burstable VM,
and only then by falling back to on-demand capacity.

The rescue capacity Burst-HADS can call on at the instant of an interruption
is, however, fixed in advance. Its burstable fleet is sized once, after the
primary scheduling step, as a fixed fraction of the instances the search
selected, and a burstable instance holds only one task at a time so that
credits continue to accrue. Neither quantity depends on how much work a
hibernation actually displaces or on how much deadline slack remains when it
occurs. When the displaced work exceeds the standing burstable capacity, every
remaining task falls through to on-demand instances, which are the most
expensive resource in the catalogue and carry a start-up latency.

This paper makes the following contributions.

\begin{itemize}
\item We present \textbf{R-BurstHADS}, which adds two provisioning behaviours to
Burst-HADS: a risk-triggered \emph{preemptive} step (Algorithm~\ref{alg:preempt})
that launches a replacement spot instance for a high-risk VM before it is
interrupted, and a \emph{deadline-aware reactive} step
(Algorithm~\ref{alg:sizing}) that provisions the minimum additional capacity
required for the displaced work to finish before the deadline. Both draw from
the same instance catalogue available to the baselines.
\item We \textbf{validate our baseline implementations} against the published
results of Burst-HADS before using them for comparison
(Section~\ref{sec:validation}).
\item We evaluate all three schedulers under the \textbf{hibernation scenarios of
the source work} across six workload sizes and 30 seeds, and show that
R-BurstHADS's cost advantage grows monotonically with the interruption rate
while its makespan advantage is essentially rate-independent
(Section~\ref{sec:results}).
\end{itemize}

\section{Background and Related Work}

\subsection{Scheduling on interruptible instances}

Work on interruptible cloud instances divides broadly into bidding and
checkpointing strategies for surviving reclamation, and scheduling strategies
that treat interruption as a first-class constraint. AutoBoT-like approaches
select instance types under a deadline and budget but do not model hibernation.
The line we build on schedules explicitly against a hibernation process.

\subsection{HADS and Burst-HADS}

HADS~\cite{hads2021} formulates BoT scheduling on spot and on-demand instances
as a cost-minimisation problem under a deadline $D$, building a static map with
a greedy single-pass construction and migrating tasks away from hibernated
instances at runtime. Burst-HADS~\cite{bursthads} replaces the greedy
construction with an Iterated Local Search (ILS), adds burstable instances as a
rescue tier, and introduces a work-stealing procedure that rebalances load
whenever an instance becomes idle.

Because Burst-HADS's Burst Migration Procedure draws burstable candidates from
the idle set and moves each recipient to the busy set on assignment, a burstable
instance accepts exactly one migrated task at a time. Combined with a fleet size
fixed at scheduling time, this bounds the burstable rescue capacity independently
of the size of the interruption --- the gap this work addresses.

\section{System Model and Problem Formulation}

An application is a set $T$ of $|T|$ independent tasks with a common deadline
$D$. A task $t_i$ has execution time $e_i$ on a reference core and memory
requirement $m_i$. The provider offers a catalogue $M$ of VM types; $vm_j$ has
processing speed $s_j$, virtual core count $\mathit{nvc}_j$, memory $M_j$ and
price $c_j$ per unit time. A multi-core instance executes up to $\mathit{nvc}_j$
tasks simultaneously, one per core, so its whole-instance throughput is
$s_j\cdot \mathit{nvc}_j$. Symbols are collected in Table~\ref{tab:vars}.

\begin{table}[t]
\caption{Variables and parameters used in the algorithms}
\label{tab:vars}
\centering
\footnotesize
\begin{tabular}{ll}
\toprule
Symbol & Meaning \\
\midrule
$T$, $t_i$ & task set; task $i$ \\
$e_i$, $m_i$ & execution time and memory requirement of $t_i$ \\
$e_{ij}$ & execution time of $t_i$ on $vm_j$ \\
$M$, $vm_j$ & VM catalogue; instance $j$ \\
$s_j$, $\mathit{nvc}_j$ & speed and virtual-core count of $vm_j$ \\
$M_j$, $c_j$ & memory capacity and price per period of $vm_j$ \\
$\Theta(\cdot)$ & throughput of a VM set, $\sum_j s_j\!\cdot\!\mathit{nvc}_j$ \\
$D$ & application deadline \\
$D_{\mathit{spot}}$ & deadline minus the worst-case migration reserve \\
$\mathit{RD}_{\mathit{spot}}$ & relaxed $D_{\mathit{spot}}$ used inside the ILS \\
$\omega$ & checkpoint overhead fraction \\
$T_{\mathit{start}}$ & instance start-up latency \\
$\mathit{cc}_j$, $\mathit{rcc}_{ij}$ & accrued and required CPU credits \\
$\lambda_h$, $\lambda_r$ & hibernation and resumption rates \\
$k_h$, $k_r$ & expected hibernations / resumptions over $D$ \\
$S$, $S_{\mathit{best}}$ & candidate and incumbent scheduling solutions \\
$\alpha$ & cost/makespan weight in the fitness function \\
$\mathit{burst\_rate}$ & burstable fleet size as a fraction of $|S|$ \\
$Q_l$ & unfinished tasks of a hibernated instance $vm_l$ \\
$\mathit{IR}$, $\mathit{BR}$ & sets of idle and busy instances \\
$M^{o}$, $M^{b}$ & elastic on-demand and burstable pools \\
$A$ & active replacement instances provisioned at runtime \\
$W$ & remaining work of the tasks awaiting placement \\
$\rho$ & preemptive provisioning risk threshold \\
$p_v$ & hibernation probability of $vm_v$ over one task \\
\bottomrule
\end{tabular}
\end{table}

Cost is instance-seconds: an instance is charged from launch until termination
or hibernation, whether or not it is executing. The objective is
\begin{equation}
\min\ \alpha \cdot \widehat{\mathit{cost}} + (1-\alpha)\cdot \widehat{\mathit{mkp}}
\quad \text{s.t.}\quad \mathit{mkp} \leq D
\label{eq:fitness}
\end{equation}
over normalised cost and makespan. Candidate instances are ranked by
\begin{equation}
\mathit{weight}(vm_j) = \frac{\mathit{Gflops}_j}{c_j} = \frac{s_j \cdot \mathit{nvc}_j}{c_j}
\label{eq:wrr}
\end{equation}
where $\mathit{Gflops}_j$ quantifies the whole instance's computing power.
Placement on spot instances is constrained not by $D$ but by
\begin{equation}
D_{\mathit{spot}} = D - \max_i \frac{e_i \cdot (1+\omega)}{s_{\min}}
\label{eq:dspot}
\end{equation}
a reserve guaranteeing that the longest task can still be migrated to the
slowest instance and executed before $D$ if its host hibernates.

\section{R-BurstHADS}
\label{sec:design}

R-BurstHADS inherits Burst-HADS's ILS (Algorithm~\ref{alg:primary}), greedy
construction (Algorithm~\ref{alg:greedy}), local search
(Algorithm~\ref{alg:ls}), burst migration (Algorithm~\ref{alg:migrate}) and
work stealing (Algorithm~\ref{alg:steal}), and adds two provisioning
behaviours, Algorithms~\ref{alg:preempt} and~\ref{alg:sizing}. Both launch
instances of the most cost-efficient spot type present in the scheduler's own
pool, $\arg\max_j s_j\mathit{nvc}_j/c_j$, so the proposed method has access to
no instance type the baselines cannot also select.

\begin{algorithm}[t]
\caption{R-BurstHADS primary scheduler}
\label{alg:primary}
\begin{algorithmic}[1]
\STATE $D_{\mathit{spot}} \leftarrow$ Eq.~(\ref{eq:dspot})
\STATE $S \leftarrow \textsc{GreedyInitial}(T, M, D_{\mathit{spot}})$
\STATE $S_{\mathit{best}} \leftarrow \textsc{LocalSearch}(S, D_{\mathit{spot}})$
\STATE $\mathit{RD}_{\mathit{spot}} \leftarrow D_{\mathit{spot}}$;\ \ $it_b \leftarrow 0$;\ \ $it_r \leftarrow 0$
\FOR{$it = 1$ to $\mathit{max\_iteration}$}
  \IF{unused spot instances remain}
    \STATE add one at random to $S.\mathit{selected}$
  \ENDIF
  \IF{$it - it_b > \mathit{max\_failed}$ \AND $it - it_r > \mathit{max\_failed}$}
    \STATE $\mathit{RD}_{\mathit{spot}} \leftarrow \min\!\big(\mathit{RD}_{\mathit{spot}}(1+\mathit{relaxed\_rate}),\, D\big)$
    \STATE $it_r \leftarrow it$
  \ENDIF
  \STATE $S' \leftarrow \textsc{LocalSearch}(S, \mathit{RD}_{\mathit{spot}})$
  \IF{$\mathit{fitness}(S') < \mathit{fitness}(S_{\mathit{best}})$}
    \STATE $S_{\mathit{best}} \leftarrow S'$;\ \ $it_b \leftarrow it$
  \ENDIF
  \STATE $S \leftarrow S'$
\ENDFOR
\STATE $n_b \leftarrow \lceil \mathit{burst\_rate}\cdot|S_{\mathit{best}}.\mathit{selected}|\rceil$
\STATE $B \leftarrow \textsc{LaunchBurstable}(n_b)$ from $M^{b}$
\FORALL{$t_i$ with predicted finish $> D_{\mathit{spot}}$, latest first}
  \STATE $b \leftarrow$ next free instance of $B$
  \IF{$b \neq \varnothing$ \AND baseline finish of $t_i$ on $b$ improves on its current finish}
    \STATE move $t_i$ to $b$ in baseline mode
  \ENDIF
\ENDFOR
\STATE \textsc{PreemptiveProvision}($S_{\mathit{best}}$) \COMMENT{Alg.~\ref{alg:preempt}}
\RETURN $S_{\mathit{best}}$
\end{algorithmic}
\end{algorithm}

\begin{algorithm}[t]
\caption{\textsc{GreedyInitial}: initial solution}
\label{alg:greedy}
\begin{algorithmic}[1]
\STATE sort $T$ by memory requirement, descending
\FORALL{$t_i \in T$}
  \FORALL{$vm_j \in \mathit{selected}$, least loaded first}
    \IF{$\textsc{Feasible}(t_i, vm_j, D_{\mathit{spot}})$}
      \STATE assign $t_i$ to $vm_j$; \textbf{continue} with next task
    \ENDIF
  \ENDFOR
  \STATE $K \leftarrow \{vm_j \in M^{s}\setminus \mathit{selected} : \textsc{Feasible}(t_i,vm_j,D_{\mathit{spot}})\}$
  \STATE $vm_j \leftarrow \textsc{WrrPick}(K)$ using Eq.~(\ref{eq:wrr})
  \IF{$vm_j \neq \varnothing$}
    \STATE select and assign; \textbf{continue} with next task
  \ENDIF
  \FORALL{$vm_j \in M^{o}$ already running, cheapest first}
    \IF{$\textsc{Feasible}(t_i, vm_j, D)$}
      \STATE assign; \textbf{continue} with next task
    \ENDIF
  \ENDFOR
  \STATE $vm_{\mathit{new}} \leftarrow \textsc{LaunchOnDemand}()$ from $M^{o}$
  \IF{$\textsc{Feasible}(t_i, vm_{\mathit{new}}, D)$}
    \STATE assign
  \ELSE
    \STATE \textbf{abort}: no feasible placement before $D$
  \ENDIF
\ENDFOR
\end{algorithmic}
\end{algorithm}

\begin{algorithm}[t]
\caption{\textsc{LocalSearch}}
\label{alg:ls}
\begin{algorithmic}[1]
\REQUIRE solution $S$, bound $d$
\STATE $\mathit{best} \leftarrow S$
\STATE $\mathit{dest} \leftarrow$ one instance chosen at random from $S.\mathit{selected}\cup$ unused spot
\FOR{$a = 1$ to $\mathit{max\_attempt}$}
  \STATE $\mathit{cand} \leftarrow \mathit{clone}(\mathit{best})$
  \STATE $R \leftarrow$ random sample of $\mathit{swap\_rate}\cdot|T|$ tasks
  \FORALL{$t_i \in R$}
    \IF{memory of $\mathit{dest}$ permits}
      \STATE move $t_i$ to $\mathit{dest}$
    \ELSE
      \STATE revert $t_i$
    \ENDIF
  \ENDFOR
  \IF{$\mathit{fitness}(\mathit{cand}, d) < \mathit{fitness}(\mathit{best}, d)$}
    \STATE $\mathit{best} \leftarrow \mathit{cand}$
  \ENDIF
\ENDFOR
\RETURN $\mathit{best}$
\end{algorithmic}
\end{algorithm}

\subsection{Preemptive replacement provisioning}

Let $p_v$ be the probability that spot instance $v$ hibernates during the
execution of an average task, derived from its rate $\lambda_h$. For each $v$
with $p_v$ above a threshold $\rho$, R-BurstHADS compares the expected cost of
rescuing $v$'s tasks reactively against the cost of holding a replacement
ready, and provisions only when the comparison favours it
(Algorithm~\ref{alg:preempt}).

\begin{algorithm}[t]
\caption{\textsc{PreemptiveProvision} (Theorem 1)}
\label{alg:preempt}
\begin{algorithmic}[1]
\REQUIRE incumbent $S_{\mathit{best}}$, threshold $\rho$
\STATE $\bar{e} \leftarrow$ mean execution time over $T$
\STATE $r_s, s_s \leftarrow$ price and speed of $\arg\max_j s_j\mathit{nvc}_j/c_j$ over $M^{s}$
\FORALL{$vm_v \in M^{s}$ with $\lambda_h(v) > 0$}
  \STATE $p_v \leftarrow 1 - e^{-\lambda_h(v)\,\bar{e}/s_v}$
  \IF{$p_v < \rho$} \STATE \textbf{continue} \ENDIF
  \STATE $n_v \leftarrow |\{t_i : S_{\mathit{best}}(t_i) = v\}|$
  \STATE $C_{\mathit{react}} \leftarrow n_v\,\bar{e}\,\big(r_{o}/s_{o} - r_s/s_s\big)$
  \STATE $C_{\mathit{proact}} \leftarrow T_{\mathit{start}}\,r_s + (\bar{e}/s_s)\,r_s$
  \IF{$p_v \cdot C_{\mathit{react}} > C_{\mathit{proact}}$}
    \STATE launch one replacement, ready at $T_{\mathit{start}}$
  \ENDIF
\ENDFOR
\end{algorithmic}
\end{algorithm}

\subsection{Deadline-aware capacity sizing}

When spot capacity is lost, R-BurstHADS sizes the replacement fleet from the
deadline rather than from a fixed queue-length target. With $W$ the remaining
work of the tasks awaiting placement, $t$ the current time and $\Theta(\cdot)$
the throughput of an instance set,
\begin{equation}
T = D - t - T_{\mathit{start}}, \quad
n_{\mathit{extra}} = \left\lceil \frac{W/T - \Theta(A)}{\theta} \right\rceil
\label{eq:sizing}
\end{equation}
where $A$ is the set of active replacements and $\theta$ the throughput of one
new instance. If the surviving fleet can already retire the queue before $D$
the deficit is non-positive and nothing is provisioned; the scheduler then
behaves exactly as Burst-HADS. Should the estimate prove optimistic, the
per-task placement path adds capacity incrementally under its own deadline
test, so no safety margin is applied to (\ref{eq:sizing}).

\begin{algorithm}[t]
\caption{\textsc{DeadlineAwareProvision} (Theorem 2)}
\label{alg:sizing}
\begin{algorithmic}[1]
\REQUIRE queued tasks $Q$, current time $t$, active replacements $A$
\STATE $W \leftarrow \sum_{t_i \in Q}$ remaining time of $t_i$
\STATE $T \leftarrow D - t - T_{\mathit{start}}$
\IF{$W \leq 0$ \OR $T \leq 0$} \RETURN $0$ \ENDIF
\STATE $\theta \leftarrow s_{\mathit{new}}\cdot \mathit{nvc}_{\mathit{new}}$
\STATE $\mathit{deficit} \leftarrow W/T - \Theta(A)$
\IF{$\mathit{deficit} \leq 0$} \RETURN $0$ \ENDIF
\STATE $n \leftarrow \min\!\big(\mathit{max\_per\_event},\ \lceil \mathit{deficit}/\theta \rceil\big)$
\STATE launch $n$ instances, ready at $t + T_{\mathit{start}}$
\STATE redistribute $Q$ over $A \cup \mathit{new}$ by earliest estimated finish
\RETURN $n$
\end{algorithmic}
\end{algorithm}

\begin{algorithm}[t]
\caption{Burst migration procedure}
\label{alg:migrate}
\begin{algorithmic}[1]
\REQUIRE $Q_l$, $\mathit{IR}$, $\mathit{BR}$, $M^{o}$, $D$
\STATE $Q_l \leftarrow \textsc{SortTasks}(Q_l)$ \COMMENT{checkpointed tasks first}
\FORALL{$t_i \in Q_l$}
  \STATE \textbf{Attempt 0} (R-BurstHADS): a provisioned replacement in $A$ with
         free capacity that finishes $t_i$ before $D$
  \STATE \textbf{Attempt 1}: burstable $vm_j \in \mathit{IR}$ with
         $\mathit{cc}_j > \mathit{rcc}_{ij}$; migrate in burst mode,
         $\mathit{IR}\!\leftarrow\!\mathit{IR}\setminus\{vm_j\}$,
         $\mathit{BR}\!\leftarrow\!\mathit{BR}\cup\{vm_j\}$
  \STATE \textbf{Attempt 2}: non-burstable $vm_j \in \textsc{SortByMarket}(\mathit{IR}\cup \mathit{BR})$
         passing $\textsc{CheckMigration}(t_i,vm_j,D)$
  \STATE \textbf{Attempt 3} (R-BurstHADS): \textsc{DeadlineAwareProvision}
         \COMMENT{Alg.~\ref{alg:sizing}}
  \STATE \textbf{Attempt 4}: new $vm_j \in M^{o}$, cheapest first, if
         $\mathit{start}_{ij} + e_{ij} + \omega < D$
\ENDFOR
\end{algorithmic}
\end{algorithm}

\begin{algorithm}[t]
\caption{Work stealing}
\label{alg:steal}
\begin{algorithmic}[1]
\REQUIRE newly idle $vm_k$, busy set $\mathit{BR}$, deadline $D$
\FORALL{non-burstable $vm_j \in \mathit{BR}$, most loaded first}
  \FORALL{waiting $t_i$ on $vm_j$}
    \IF{$\textsc{Slack}(t_i, vm_k, D) > 0$}
      \STATE move $t_i$ to $vm_k$
      \IF{$vm_k$ is burstable}
        \STATE mark $t_i$ baseline mode; \RETURN \COMMENT{one task only}
      \ENDIF
    \ENDIF
  \ENDFOR
\ENDFOR
\end{algorithmic}
\end{algorithm}

\section{Experimental Setup}
\label{sec:setup}

\subsection{Simulator, catalogue and workload}

We use a discrete-event simulator in which each instance executes up to
$\mathit{nvc}_j$ tasks concurrently. Hibernation stops the charge and
checkpoints resident tasks; a resumed instance continues from its checkpoint.
The shared pool (Table~\ref{tab:vms}) contains three instances of each spot
type, one burstable instance seeding the elastic pool $M^{b}$, and one
on-demand instance as the template for the unbounded pool $M^{o}$. Prices are
us-east-1 rates. Per-instance-type hibernation rates use the source work's
worst-case parameterisation ($k_h=5$ over $D=2700$\,s) for \texttt{m5.xlarge}
and a low nominal rate for the remaining spot types; this differentiation is a
modelling choice of the present study rather than a measured quantity.

Tasks are drawn independently with execution time
$\mathcal{U}[100,350]$\,s and memory $\mathcal{U}[2.85,13.19]$\,MB, matching the
ranges reported for the reference jobs, with
$|T| \in \{10,20,50,100,200,300\}$.

\begin{table}[t]
\caption{VM catalogue. Three instances of each spot type.}
\label{tab:vms}
\centering
\footnotesize
\begin{tabular}{llccccc}
\toprule
Type & Market & $s_j$ & $\mathit{nvc}_j$ & Mem & \$/h & $\lambda_h$ \\
\midrule
c5.large  & spot      & 2 & 2 & 4\,GB  & 0.0306 & $0.03/3600$ \\
c5.xlarge & spot      & 4 & 4 & 8\,GB  & 0.0612 & $0.04/3600$ \\
m5.xlarge & spot      & 4 & 4 & 16\,GB & 0.0700 & $5.0/2700$ \\
t3.large  & burstable & 2 & 1 & 8\,GB  & 0.0832 & --- \\
c5.large  & on-demand & 2 & 2 & 4\,GB  & 0.0850 & --- \\
\bottomrule
\end{tabular}
\end{table}

\subsection{Deadline derivation}

The deadline is derived from the pool rather than fixed, so that experimental
difficulty is invariant to pool size:
\begin{equation}
M^{*} = \max\!\left(\frac{|T|\bar{e}}{\Theta(M^{s})},\ \frac{e_{\max}}{s_{\max}}\right),
\qquad D = \mathit{DF}\cdot\kappa\cdot M^{*}
\label{eq:deadline}
\end{equation}
with $\kappa = 3$ calibrated so that $\mathit{DF}=1$ is met by a well-packed
undisturbed run but not by a disrupted one without escalation. A floor keeps
$D$ above the migration reserve of (\ref{eq:dspot}), below which the problem is
infeasible rather than tight. All reported experiments use $\mathit{DF}=1$.

\subsection{Hibernation scenarios and parameters}

Interruptions follow the source work's scenarios
(Table~\ref{tab:scenarios}), in which each spot instance is subject to an
independent Poisson process with rate $\lambda_h = k_h/D$ and, where resumption
is modelled, $\lambda_r = k_r/D$. Algorithm parameters are given in
Table~\ref{tab:params}. Every configuration is run with 30 seeds; all three
schedulers see a byte-identical task set and interruption stream for a given
seed, so the comparison is paired at the level of the workload. In total the
evaluation comprises 2{,}700 simulation runs.

\begin{table}[t]
\caption{Hibernation scenarios}
\label{tab:scenarios}
\centering
\footnotesize
\begin{tabular}{lccccc}
\toprule
 & sc1 & sc2 & sc3 & sc4 & sc5 \\
\midrule
$k_h$ & 1 & 5 & 1 & 5 & 3 \\
$k_r$ & 0 & 0 & 5 & 5 & 2.5 \\
\bottomrule
\end{tabular}
\end{table}

\begin{table}[t]
\caption{Algorithm and experiment parameters}
\label{tab:params}
\centering
\footnotesize
\begin{tabular}{ll}
\toprule
Parameter & Value \\
\midrule
$\alpha$ & 0.5 \\
$\mathit{max\_iteration}$, $\mathit{max\_attempt}$ & 200, 50 \\
$\mathit{swap\_rate}$, $\mathit{max\_failed}$ & 0.10, 20 \\
$\mathit{relaxed\_rate}$ & 0.25 \\
$\mathit{burst\_rate}$ & 0.2 \\
Allocation cycle & 900\,s \\
$T_{\mathit{start}}$, $\omega$ & 45\,s, 0.10 \\
$\rho$, $\mathit{max\_per\_event}$ & 0.05, 10 \\
$\mathit{DF}$, $\kappa$ & 1.0, 3 \\
Seeds per configuration & 30 \\
\bottomrule
\end{tabular}
\end{table}

\section{Baseline Validation}
\label{sec:validation}

A comparison is only as trustworthy as its baseline implementation. We
therefore first reproduce the published comparison between HADS and Burst-HADS,
using the source work's own catalogue, job definitions, fixed deadline
$D=2700$\,s and scenarios. Table~\ref{tab:validation} and
Fig.~\ref{fig:validation} report the outcome.

\begin{table}[t]
\caption{Baseline validation against published results}
\label{tab:validation}
\centering
\footnotesize
\begin{tabular}{lccccc}
\toprule
 & \multicolumn{2}{c}{HADS mkp (s)} & \multicolumn{2}{c}{$\Delta$ cost (\%)} & $\Delta$ mkp \\
\cmidrule(lr){2-3}\cmidrule(lr){4-5}
Job & ours & pub. & ours & pub. & ours \\
\midrule
J60   & 2476 & 2620 & $+29.5$ & $+30.8$ & $-33.9$ \\
J80   & 2501 & 2581 & $+25.0$ & $+11.3$ & $-36.0$ \\
J100  & 2541 & 2518 & $+31.3$ & --- & $-8.4$ \\
ED200 & 2694 & 2680 & $+21.6$ & $+23.6$ & $-2.7$ \\
\bottomrule
\end{tabular}
\end{table}

\begin{figure}[t]
\centering
\includegraphics[width=\columnwidth]{fig/fig10_validation.png}
\caption{Validation against the published results. Our HADS lands within a few
percent of the reported makespans, pinned just below the deadline.}
\label{fig:validation}
\end{figure}

Our HADS implementation lands within a few percent of the published makespans,
in every case pinned just below the deadline, which is the expected behaviour of
a cost-minimising greedy that packs to its budget. The direction of the cost
difference is reproduced in every job and its magnitude closely in three of
four: Burst-HADS is faster than HADS and slightly more expensive, which is what
the source work reports.

Our reproduction of the makespan \emph{reduction} is smaller than published, and
the shortfall grows with job size. The most plausible explanation is capacity: at
a fixed $D=2700$\,s the largest job consumes the whole pool, leaving the search
no spare instances to spread onto. We report this as a limitation of the
reproduction rather than a correction to the source work.

\section{Results}
\label{sec:results}

\begin{figure*}[t]
\centering
\includegraphics[width=\textwidth]{fig/fig1_makespan_cost.png}
\caption{Makespan (top) and monetary cost (bottom) against workload size for
each hibernation scenario. Means over 30 seeds. Lower is better.}
\label{fig:main}
\end{figure*}

\begin{figure*}[t]
\centering
\includegraphics[width=\textwidth]{fig/fig6_makespan_variability.png}
\caption{Makespan with $\pm1$ standard deviation over the 30 seeds. Dispersion
grows with the interruption rate but the ordering is stable.}
\label{fig:variability}
\end{figure*}

\subsection{Task completion and deadline satisfaction}

All three schedulers completed 100\% of tasks in all 30 configurations, so
completion rate does not discriminate between them. Deadline satisfaction does.
Burst-HADS and R-BurstHADS met the deadline in every configuration. HADS missed
in three --- sc1, sc2 and sc3 at $|T|=50$ --- the workload size at which the
deadline is closest to its feasibility floor and HADS, lacking burstable rescue
capacity, must absorb an interruption entirely on on-demand instances subject to
start-up latency.

\subsection{Burst-HADS against HADS}

At $|T|=300$ the ordering reported in the source work is reproduced in all five
scenarios: Burst-HADS is faster than HADS, by between 1.1\% (sc2) and 24.9\%
(sc3), and more expensive, by between 1.7\% (sc4) and 16.1\% (sc1). The ordering
weakens as the workload shrinks --- four of five scenarios at $|T|=200$, two of
five at $|T|=100$ --- because a smaller workload leaves the pool underutilised
and the two schedulers converge on similar placements. We therefore treat
$|T| \geq 200$ as the regime in which this comparison is informative.
Fig.~\ref{fig:speedup} shows the speedup of both extensions relative to HADS,
and Fig.~\ref{fig:util} the fraction of the deadline each scheduler consumes.

\begin{figure*}[t]
\centering
\includegraphics[width=\textwidth]{fig/fig4_speedup.png}
\caption{Speedup relative to HADS. Values above unity indicate a faster
scheduler.}
\label{fig:speedup}
\end{figure*}

\begin{figure*}[t]
\centering
\includegraphics[width=\textwidth]{fig/fig9_deadline_utilisation.png}
\caption{Deadline utilisation, i.e. makespan divided by $D$. HADS consumes its
budget almost entirely; R-BurstHADS retains headroom.}
\label{fig:util}
\end{figure*}

\subsection{R-BurstHADS against Burst-HADS}

Fig.~\ref{fig:main} shows the full comparison, Fig.~\ref{fig:advantage}
summarises R-BurstHADS's advantage, and Fig.~\ref{fig:heatmap} resolves it per
cell. Two patterns are visible.

First, the \textbf{makespan advantage is large and rate-independent}:
R-BurstHADS is 26.0--39.8\% faster than Burst-HADS in every scenario, with no
systematic relationship to $k_h$. This is the expected signature of elastic
provisioning: whenever spot capacity is lost, replacement capacity is acquired
in proportion to the displaced work rather than drawn from a standing pool of
fixed size.

Second, the \textbf{cost advantage scales monotonically with the interruption
rate} (Fig.~\ref{fig:kh}). At $k_h=1$ the two schedulers are indistinguishable
on cost ($+0.5\%$ in sc1, $-0.5\%$ in sc3); at $k_h=3$ R-BurstHADS is 12.5\%
cheaper; at $k_h=5$ it is 14.3\% cheaper when instances resume (sc4) and 64.3\%
cheaper when they do not (sc2). Under frequent interruption Burst-HADS exhausts
its standing burstable capacity early and places the remainder of the displaced
work on on-demand instances, the most expensive resource available, whereas
R-BurstHADS replaces lost spot capacity with spot capacity. That the cost
advantage tracks a parameter which was not tuned, and which controls the
environment rather than the algorithm, is the strongest evidence that the effect
is mechanistic rather than an artefact of configuration.

Fig.~\ref{fig:tradeoff} places both extensions in the cost--makespan plane
relative to HADS, and Fig.~\ref{fig:cpt} normalises cost per task.
Table~\ref{tab:advantage} gives the summary figures and Table~\ref{tab:main}
the full grid.

\begin{figure}[t]
\centering
\includegraphics[width=\columnwidth]{fig/fig2_rburst_advantage.png}
\caption{R-BurstHADS's advantage over Burst-HADS, averaged over $|T|\geq100$.}
\label{fig:advantage}
\end{figure}

\begin{figure}[t]
\centering
\includegraphics[width=\columnwidth]{fig/fig8_kh_sensitivity.png}
\caption{Advantage against the expected interruption count $k_h$. The cost
advantage is monotone in $k_h$; the makespan advantage is not.}
\label{fig:kh}
\end{figure}

\begin{figure}[t]
\centering
\includegraphics[width=\columnwidth]{fig/fig5_tradeoff.png}
\caption{Cost--makespan plane relative to HADS, for $|T|\geq100$. R-BurstHADS
occupies the faster region; Burst-HADS clusters near the baseline.}
\label{fig:tradeoff}
\end{figure}

\begin{figure*}[t]
\centering
\includegraphics[width=\textwidth]{fig/fig7_heatmap.png}
\caption{Per-cell advantage of R-BurstHADS over Burst-HADS.}
\label{fig:heatmap}
\end{figure*}

\begin{figure*}[t]
\centering
\includegraphics[width=\textwidth]{fig/fig3_cost_per_task.png}
\caption{Cost per task. Normalising by workload size isolates the per-unit
efficiency of each scheduler.}
\label{fig:cpt}
\end{figure*}

\begin{table}[t]
\caption{R-BurstHADS relative to Burst-HADS, mean over $|T|\geq100$}
\label{tab:advantage}
\centering
\footnotesize
\begin{tabular}{lccrr}
\toprule
Scenario & $k_h$ & $k_r$ & Cost (\%) & Makespan (\%) \\
\midrule
sc1 & 1 & 0   & $+0.5$  & $+39.8$ \\
sc3 & 1 & 5   & $-0.5$  & $+35.6$ \\
sc5 & 3 & 2.5 & $+12.5$ & $+29.7$ \\
sc4 & 5 & 5   & $+14.3$ & $+28.3$ \\
sc2 & 5 & 0   & $+64.3$ & $+26.0$ \\
\bottomrule
\end{tabular}
\end{table}

\begin{table*}[t]
\caption{Full results. Monetary cost (USD) and makespan (s), means over 30 seeds.}
\label{tab:main}
\centering
\begin{tabular}{llcccccc}
\toprule
& & \multicolumn{2}{c}{HADS} & \multicolumn{2}{c}{Burst-HADS} & \multicolumn{2}{c}{R-BurstHADS} \\
\cmidrule(lr){3-4}\cmidrule(lr){5-6}\cmidrule(lr){7-8}
Sc. & $n$ & cost & mkp & cost & mkp & cost & mkp \\
\midrule
\multirow{6}{*}{sc1} & 10 & 0.0400 & 114 & 0.0321 & 108 & 0.0317 & 107 \\
 & 20 & 0.0658 & 128 & 0.0530 & 125 & 0.0530 & 125 \\
 & 50 & 0.1150 & 295 & 0.0868 & 208 & 0.0831 & 198 \\
 & 100 & 0.1160 & 533 & 0.1125 & 525 & 0.1230 & 256 \\
 & 200 & 0.1670 & 1177 & 0.1938 & 1121 & 0.1934 & 698 \\
 & 300 & 0.2141 & 1836 & 0.2486 & 1449 & 0.2356 & 908 \\
\midrule
\multirow{6}{*}{sc2} & 10 & 0.0272 & 150 & 0.0190 & 137 & 0.0190 & 137 \\
 & 20 & 0.0562 & 223 & 0.0388 & 192 & 0.0348 & 185 \\
 & 50 & 0.1947 & 316 & 0.2043 & 300 & 0.0954 & 279 \\
 & 100 & 0.2762 & 624 & 0.2643 & 614 & 0.0753 & 361 \\
 & 200 & 0.3950 & 1248 & 0.4017 & 1234 & 0.1563 & 1065 \\
 & 300 & 0.5053 & 1874 & 0.5292 & 1852 & 0.1945 & 1310 \\
\midrule
\multirow{6}{*}{sc3} & 10 & 0.1003 & 112 & 0.0937 & 119 & 0.0937 & 119 \\
 & 20 & 0.1234 & 124 & 0.1128 & 125 & 0.1128 & 125 \\
 & 50 & 0.1731 & 204 & 0.1443 & 201 & 0.1416 & 194 \\
 & 100 & 0.1707 & 487 & 0.1771 & 528 & 0.1848 & 254 \\
 & 200 & 0.2248 & 1124 & 0.2570 & 1067 & 0.2592 & 701 \\
 & 300 & 0.2709 & 1759 & 0.3072 & 1320 & 0.3007 & 923 \\
\midrule
\multirow{6}{*}{sc4} & 10 & 0.0953 & 126 & 0.0813 & 136 & 0.0813 & 136 \\
 & 20 & 0.1044 & 151 & 0.0894 & 168 & 0.0889 & 167 \\
 & 50 & 0.1684 & 234 & 0.1719 & 256 & 0.1393 & 244 \\
 & 100 & 0.1884 & 499 & 0.1854 & 548 & 0.1503 & 297 \\
 & 200 & 0.2615 & 1080 & 0.2713 & 1123 & 0.2417 & 928 \\
 & 300 & 0.3364 & 1696 & 0.3421 & 1552 & 0.2929 & 1087 \\
\midrule
\multirow{6}{*}{sc5} & 10 & 0.0880 & 122 & 0.0762 & 133 & 0.0762 & 133 \\
 & 20 & 0.0985 & 142 & 0.0873 & 157 & 0.0864 & 151 \\
 & 50 & 0.1689 & 235 & 0.1700 & 260 & 0.1273 & 245 \\
 & 100 & 0.1781 & 519 & 0.1715 & 553 & 0.1451 & 295 \\
 & 200 & 0.2480 & 1135 & 0.2511 & 1102 & 0.2313 & 903 \\
 & 300 & 0.3165 & 1758 & 0.3255 & 1585 & 0.2780 & 1079 \\
\bottomrule
\end{tabular}
\end{table*}

\section{Discussion}

\subsection{When the contribution is inert}

R-BurstHADS provisions nothing when the deadline is not threatened, which is a
property of (\ref{eq:sizing}) rather than an accident. In the validation
experiments, where $D$ is fixed at 2700\,s against a pool sized for a much
shorter makespan, R-BurstHADS is identical to Burst-HADS in most cells. The
exception is the largest job, where memory pressure binds and it is 33\%
cheaper. This is the intended behaviour: the mechanism costs nothing when it is
not needed, and its value is confined to deadline-pressured regimes.

\subsection{Cost--makespan trade}

Neither extension dominates HADS on cost. HADS is the cheapest of the three in
most configurations, because it packs onto the fewest instances and consumes its
deadline budget --- visible directly in Fig.~\ref{fig:util}. What the extensions
buy is time and deadline safety: HADS is the only scheduler in the study that
misses a deadline, and it is 26--40\% slower than R-BurstHADS at every workload
size above 100 tasks. Which of these is preferable is a property of the
application, not of the scheduler, and we report both rather than collapsing
them into a single figure of merit.

\subsection{Limitations}

The evaluation is a simulation and inherits the abstractions of its model. Task
execution times are drawn from a uniform distribution whose range matches the
reference jobs but whose shape does not reflect the right-skewed,
straggler-heavy profile of real workloads; this understates the tail that
scheduling must absorb. Hibernation stops compute billing but incurs no storage
charge, where the provider does charge for the retained EBS volume; the omitted
quantity is small relative to compute but non-zero. The per-instance-type
differentiation of hibernation rates, as noted in Section~\ref{sec:setup}, is a
modelling choice of this study. Finally, the baseline reproduction degrades on
the largest job, so conclusions drawn at that scale rest on a weaker validation
than those at smaller ones.

\section{Conclusion}

We presented R-BurstHADS, which extends Burst-HADS with elastic spot
provisioning sized against the remaining deadline by a work-conservation
argument, drawing from the same instance catalogue available to the baselines.
Across 2{,}700 simulation runs it reduced makespan relative to Burst-HADS by
26.0--39.8\% in every hibernation scenario, and its monetary-cost advantage
scaled monotonically with the expected interruption count, from cost-neutral at
$k_h=1$ to 64.3\% cheaper at $k_h=5$ without resumption. Both extensions met the
deadline in every configuration, where HADS did not.

Two directions follow. The provisioning rule currently treats replacement
instances as equally reliable as those they replace; incorporating the
replacement's own interruption risk into (\ref{eq:sizing}) would make the sizing
risk-aware as well as deadline-aware. Second, the evaluation should move from a
synthetic uniform workload to real BoT traces, where the straggler tail is the
dominant source of makespan and the value of elastic rescue capacity is likely
to be larger than reported here.

\begin{thebibliography}{00}
\bibitem{hads2021} L. Teylo, L. Arantes, P. Sens, and L. M. de A. Drummond,
``A dynamic task scheduler tolerant to multiple hibernations in cloud
environments,'' \emph{Cluster Computing}, vol.~24, no.~2, pp.~1051--1073, 2021.

\bibitem{bursthads} L. Teylo, L. Arantes, P. Sens, and L. M. de A. Drummond,
``Scheduling bag-of-tasks in clouds using spot and burstable virtual
machines,'' \emph{IEEE Transactions on Cloud Computing}, 2023.

\bibitem{hadsconf} L. Teylo, L. Arantes, P. Sens, and L. M. de A. Drummond,
``A hibernation aware dynamic scheduler for cloud environments,'' in
\emph{Proc. 48th Int. Conf. on Parallel Processing Workshops (ICPP)}, 2019.

\bibitem{aws} Amazon Web Services, ``Amazon EC2 Spot Instances,''
\url{https://aws.amazon.com/ec2/spot/}.

\bibitem{t3} Amazon Web Services, ``Burstable performance instances,''
\emph{Amazon EC2 User Guide}.

\end{thebibliography}

\end{document}
